\documentclass{beamer}
\hypersetup{pdfpagemode=FullScreen}
\setbeamertemplate{navigation symbols}{}
\usepackage[english]{babel}
\usepackage[latin1]{inputenc}
\usepackage{amsmath,amssymb,amsthm}
\usepackage[T1]{fontenc}
\usepackage{times}
\usepackage{mathptmx}
\usepackage{graphicx}
\usepackage{booktabs}
\usepackage{hyperref}
\usepackage{xcolor}
\usepackage{mathtools}
\usepackage{algorithm,algorithmic}
\usepackage{bm}

\renewcommand{\algorithmicrequire}{\textbf{Input:}}
\renewcommand{\algorithmicensure}{\textbf{Output:}}

\mode<presentation>{
  \usetheme{Madrid}
  \usecolortheme{whale}
}

\setbeamertemplate{blocks}[rounded][shadow=false]
\AtBeginEnvironment{block}{\small}
\AtBeginEnvironment{alertblock}{\small}
\AtBeginEnvironment{exampleblock}{\small}
\setbeamersize{text margin left=6pt, text margin right=6pt}
\setbeamercovered{transparent}
\setbeamertemplate{bibliography item}[text]
\setbeamertemplate{theorems}[numbered]
\setbeamerfont{title}{size=\large}
\setbeamerfont{date}{size=\tiny}
\setbeamertemplate{itemize items}[ball]
\setbeamertemplate{itemize subitem}[circle]
\setbeamertemplate{itemize subsubitem}[triangle]
\setbeamertemplate{caption}[numbered]

%------------------------------------------------------------------
%  TITLE PAGE
%------------------------------------------------------------------
\title[Autoencoders \& Regularisation]{%
  Autoencoders\\[4pt]
  and Regularisation}

\author[Vatsav Reddy]{%
  \texorpdfstring{%
    {\Large Vatsav Reddy}\\[6pt]
    {\small Under the Guidance of}\\[2pt]
    {\small Prof.\ Praveen Parchuri}%
  }{Vatsav Reddy}}

\institute[IIITH]{%
  {\normalsize International Institute of Information Technology, Hyderabad}}

\date{\today}

\begin{document}

\begin{frame}
  \titlepage
\end{frame}

\begin{frame}[shrink=20]{Contents}
  \tableofcontents
\end{frame}

\AtBeginSection[]{
  \begin{frame}<beamer>
    \frametitle{Current Section}
    \tableofcontents[currentsection]
  \end{frame}
}

%==================================================================
\section{Introduction to Autoencoders}
%==================================================================

\begin{frame}[shrink=10]{What is an Autoencoder?}
  \begin{block}{Definition}
    An \textbf{autoencoder} is a special feed-forward neural network that
    learns to copy its input to its output through a compressed
    intermediate representation. It performs two operations:
    \begin{itemize}\setlength\itemsep{4pt}
      \item \textbf{Encode:} map input $\mathbf{x}_i$ into a hidden
            representation $\mathbf{h}$.
      \item \textbf{Decode:} reconstruct the input from $\mathbf{h}$ to
            produce the output $\hat{\mathbf{x}}_i$.
    \end{itemize}
  \end{block}

  \begin{block}{The Two Maps}
    \[
      \mathbf{h} = g(W\mathbf{x}_i + \mathbf{b})
      \qquad\qquad
      \hat{\mathbf{x}}_i = f(W^{*}\mathbf{h} + \mathbf{c})
    \]
    with $\mathbf{x}_i \in \mathbb{R}^{n}$, $\mathbf{h} \in \mathbb{R}^{d}$,
    $W \in \mathbb{R}^{d\times n}$, $W^{*} \in \mathbb{R}^{n\times d}$.
    Here $g$ is the encoder \emph{nonlinearity} and $f$ the decoder
    nonlinearity.
  \end{block}

  \begin{alertblock}{Goal}
    We want $\mathbf{h}$ to capture the \emph{most important
    characteristics} of $\mathbf{x}_i$: everything required to reconstruct
    the original input.
    Training minimises a loss that forces
    $\hat{\mathbf{x}}_i \approx \mathbf{x}_i$.
  \end{alertblock}
\end{frame}

%------------------------------------------------------------------
\begin{frame}[shrink=10]{Undercomplete Autoencoders}
  \begin{block}{When $\dim(\mathbf{h}) < \dim(\mathbf{x}_i)$}
    The hidden layer is \emph{smaller} than the input. If the network can
    still reconstruct $\mathbf{x}_i$ from $\mathbf{h}$, then $\mathbf{h}$
    is a \textbf{loss-free encoding}: it has captured all the important
    characteristics of $\mathbf{x}_i$ in fewer dimensions.
  \end{block}

  \begin{exampleblock}{The Analogy with PCA}
    This is exactly what PCA does. The bottleneck $\mathbf{h}$:
    \begin{itemize}\setlength\itemsep{3pt}
      \item gets rid of \textbf{noise},
      \item discards \textbf{low-variance} directions,
      \item removes \textbf{correlated / redundant} dimensions.
    \end{itemize}
    We will later prove that a \emph{linear} undercomplete autoencoder is
    mathematically equivalent to PCA.
  \end{exampleblock}

  \begin{block}{Compression Intuition}
    Think of a $10$-bit object compressed to a $4$-bit code, then expanded
    back to $10$ bits. If reconstruction succeeds, the $4$-bit code retained
    all essential information --- the rest was redundant.
  \end{block}
\end{frame}

%------------------------------------------------------------------
\begin{frame}[shrink=10]{Overcomplete Autoencoders}
  \begin{block}{When $\dim(\mathbf{h}) \geq \dim(\mathbf{x}_i)$}
    The hidden layer is \emph{at least as large} as the input.
    In this case the autoencoder can learn a \textbf{trivial identity
    encoding}: simply copy $\mathbf{x}_i$ into $\mathbf{h}$, then copy
    $\mathbf{h}$ into $\hat{\mathbf{x}}_i$.
  \end{block}

  \begin{alertblock}{Why the Identity Encoding is Useless}
    An identity mapping reconstructs perfectly but tells us
    \emph{nothing} about the important characteristics of the data ---
    no compression, no insight, no denoising. It is like copying a
    file rather than compressing it.
  \end{alertblock}

  \begin{block}{The Core Problem}
    There is \emph{no reason} the network should choose to disentangle
    features rather than copy them. Left alone, it tends to learn the
    easy identity map. Hence we need \textbf{special methods}
    (regularisation) to prevent this kind of trivial representation ---
    this motivates the entire second half of these notes.
  \end{block}
\end{frame}

%==================================================================
\section{Choice of Decoder and Loss Function}
%==================================================================

\begin{frame}[shrink=10]{Two Design Choices}
  \begin{block}{What Must Be Chosen}
    Designing an autoencoder requires two decisions:
    \begin{enumerate}\setlength\itemsep{4pt}
      \item the decoder activation $f$ (and encoder activation $g$),
      \item the reconstruction \textbf{loss function}.
    \end{enumerate}
    Both depend on the \emph{type} of input data (binary vs.\ real-valued).
  \end{block}

  \begin{block}{Binary Inputs}
    Suppose each $x_{ij} \in \{0,1\}$ is a binary variable.
    Candidate decoders:
    \[
      \hat{x}_i = \tanh(W^{*}\mathbf{h}+\mathbf{c}) \in [-1,1],\quad
      \hat{x}_i = W^{*}\mathbf{h}+\mathbf{c} \in \mathbb{R},\quad
      \hat{x}_i = \text{logistic}(W^{*}\mathbf{h}+\mathbf{c})
    \]
    Since inputs lie in $\{0,1\}$, the \textbf{logistic (sigmoid)} decoder
    is most apt --- it naturally restricts outputs to $[0,1]$ and can be
    read as a probability. The encoder $g$ is also typically the sigmoid.
  \end{block}

  \begin{exampleblock}{Real-Valued Inputs}
    For $\mathbf{x}_i \in \mathbb{R}^n$ (e.g.\ pixel intensities
    $0.25,\,0.5,\ldots$), logistic and $\tanh$ would wrongly restrict the
    output to $[0,1]$ or $[-1,1]$. The most apt decoder is
    \textbf{linear}: $\hat{x}_i = W^{*}\mathbf{h}+\mathbf{c}$, since the
    output should lie in $\mathbb{R}$. Again $g$ is typically the sigmoid.
  \end{exampleblock}
\end{frame}

%------------------------------------------------------------------
\begin{frame}[shrink=10]{Loss Function: Real-Valued Inputs (Squared Error)}
  \begin{block}{Objective}
    The objective of the autoencoder is to reconstruct $\hat{\mathbf{x}}_i$
    as close to $\mathbf{x}_i$ as possible. For real-valued inputs we use
    \textbf{squared-error} loss:
    \[
      \min_{W,W^{*},\mathbf{b},\mathbf{c}}
        \frac{1}{m}\sum_{i=1}^{m}\sum_{j=1}^{n}
        \bigl(\hat{x}_{ij} - x_{ij}\bigr)^2
      \;=\;
      \min_{W,W^{*},\mathbf{b},\mathbf{c}}
        \frac{1}{m}\sum_{i=1}^{m}
        (\hat{\mathbf{x}}_i - \mathbf{x}_i)^{\!\top}
        (\hat{\mathbf{x}}_i - \mathbf{x}_i)
    \]
    This is just the average squared $\ell_2$ distance between the
    reconstruction and the input over all $m$ training points.
  \end{block}

  \begin{block}{Training via Back-Propagation}
    The autoencoder is trained exactly like a regular feed-forward network
    using back-propagation. All we need is a formula for the gradient of
    the loss with respect to the output of the last layer.
    With $L(\theta) = (\hat{\mathbf{x}}_i - \mathbf{x}_i)^{\!\top}
    (\hat{\mathbf{x}}_i - \mathbf{x}_i)$:
    \[
      \frac{\partial L(\theta)}{\partial \hat{\mathbf{x}}_i}
        = \nabla_{\hat{\mathbf{x}}_i}
          (\hat{\mathbf{x}}_i - \mathbf{x}_i)^{\!\top}
          (\hat{\mathbf{x}}_i - \mathbf{x}_i)
        = 2(\hat{\mathbf{x}}_i - \mathbf{x}_i)
    \]
    From this seed gradient, the chain rule propagates derivatives back
    through decoder and encoder, exactly as in standard back-propagation.
  \end{block}
\end{frame}

%------------------------------------------------------------------
\begin{frame}[shrink=10]{Loss Function: Binary Inputs (Cross-Entropy)}
  \begin{block}{Why a Probabilistic Loss?}
    With a sigmoid decoder the output $\sigma(\cdot)\in(0,1)$ can be read
    as a \emph{probability}: an output of $0.8$ means the network is
    $80\%$ confident the true value is $1$, and $0.2$ means $20\%$
    confident it is $1$ (i.e.\ $80\%$ confident it is $0$).
    Whenever a probabilistic interpretation is possible, \textbf{cross-entropy}
    loss often works much better than squared loss.
  \end{block}

  \begin{block}{Cross-Entropy Reconstruction Loss}
    \[
      \boxed{\;
        \min\; -\frac{1}{m}\sum_{i=1}^{m}\sum_{j=1}^{n}
          \Bigl[\,x_{ij}\log \hat{x}_{ij}
            + (1-x_{ij})\log(1-\hat{x}_{ij})\,\Bigr]
      \;}
    \]
    derived from the general cross-entropy $-\sum_l p_l\log q_l$ with the
    true distribution $p=(1-x_{ij},\,x_{ij})$ and estimate
    $q=(1-\hat{x}_{ij},\,\hat{x}_{ij})$.
  \end{block}

  \begin{exampleblock}{Which $\hat{x}_{ij}$ Minimises This?}
    Treating one term: if $x_{ij}=1$ the loss is $-\log\hat{x}_{ij}$,
    minimised by $\hat{x}_{ij}=1$. If $x_{ij}=0$ the loss is
    $-\log(1-\hat{x}_{ij})$, minimised by $\hat{x}_{ij}=0$. So the optimum
    is precisely $\hat{x}_{ij}=x_{ij}$ --- the loss rewards matching the
    target. \quad(With $h_{2j}=\hat{x}_{ij}=\sigma(a_{2j})$.)
  \end{exampleblock}
\end{frame}

%------------------------------------------------------------------
\begin{frame}[shrink=10]{Back-Propagation through the Cross-Entropy Loss}
  \begin{block}{Seed Gradient at the Output}
    For a single output unit with target $x_{ij}$ and reconstruction
    $\hat{x}_{ij}=\sigma(a_{2j})$, the loss
    $L(\theta)=-\bigl[x_{ij}\log\hat{x}_{ij}+(1-x_{ij})\log(1-\hat{x}_{ij})\bigr]$
    gives
    \[
      \frac{\partial L(\theta)}{\partial \hat{x}_{ij}}
        = -\frac{x_{ij}}{\hat{x}_{ij}}
          + \frac{1-x_{ij}}{1-\hat{x}_{ij}}
    \]
    and through the sigmoid pre-activation $a_{2j}$,
    $\dfrac{d\hat{x}_{ij}}{d a_{2j}} = \sigma(a_{2j})\bigl(1-\sigma(a_{2j})\bigr)$.
  \end{block}

  \begin{block}{Chain Rule Assembles the Full Gradient}
    \[
      \frac{\partial L(\theta)}{\partial W}
        = \underbrace{\frac{\partial L(\theta)}{\partial \mathbf{h}_2}}_{\text{reuse}}
          \,\frac{\partial \mathbf{h}_2}{\partial \mathbf{a}_2}
          \,\frac{\partial \mathbf{a}_2}{\partial \mathbf{h}_1}
          \,\frac{\partial \mathbf{h}_1}{\partial \mathbf{a}_1}
          \,\frac{\partial \mathbf{a}_1}{\partial W}
    \]
    The boxed quantity is computed once and \emph{reused};
    only the term that varies in the summation index $j$ gets
    differentiated. This is the standard back-propagation recipe applied
    to the autoencoder.
  \end{block}

  \begin{alertblock}{Summary of Choices}
    \textbf{Binary input} $\Rightarrow$ sigmoid decoder $+$ cross-entropy loss.\quad
    \textbf{Real-valued input} $\Rightarrow$ linear decoder $+$ squared-error loss.
    The encoder $g$ is typically sigmoid in both cases. The hidden layer is the
    \emph{bottleneck} that defines the learned representation.
  \end{alertblock}
\end{frame}

%==================================================================
\section{Link between PCA and Autoencoders}
%==================================================================

\begin{frame}[shrink=10]{Claim: A Linear Autoencoder is PCA}
  \begin{block}{The Four Conditions}
    The \emph{encoder} of an autoencoder is equivalent to PCA provided we:
    \begin{enumerate}\setlength\itemsep{3pt}
      \item use a \textbf{linear encoder} ($\mathbf{h}=W\mathbf{x}$, not $\sigma(W\mathbf{x})$),
      \item use a \textbf{linear decoder} ($\hat{\mathbf{x}}=W^{*}\mathbf{h}$),
      \item use a \textbf{squared-error} loss function,
      \item \textbf{normalise} the inputs as below.
    \end{enumerate}
  \end{block}

  \begin{block}{Normalisation (Centring + Scaling)}
    \[
      \hat{x}_{ij}
        = \frac{1}{\sqrt{m}}
          \Bigl(x_{ij} - \tfrac{1}{m}\sum_{k=1}^{m}x_{kj}\Bigr)
    \]
    The bracket subtracts the column (feature) mean, so the data has
    \textbf{zero mean} along each dimension $j$. Let $X'$ be this
    zero-mean data matrix; then $X = \tfrac{1}{\sqrt{m}}X'$.
  \end{block}

  \begin{exampleblock}{Why This Matters}
    With this scaling, $X^{\top}X = \frac{1}{m}(X')^{\top}X'$ is exactly the
    \textbf{covariance matrix} --- and the covariance matrix is what PCA
    diagonalises. This is the bridge between the two methods.
  \end{exampleblock}
\end{frame}

%------------------------------------------------------------------
\begin{frame}[shrink=10]{Reformulating the Objective via the Frobenius Norm}
  \begin{block}{Squared Error as a Matrix Norm}
    The reconstruction objective
    $\frac{1}{m}\sum_{i=1}^m\sum_{j=1}^n (x_{ij}-\hat{x}_{ij})^2$
    is precisely a \textbf{Frobenius norm}:
    \[
      \min_{W,W^{*}}
        \bigl\| X - H W^{*} \bigr\|_F^{2}
    \]
    where row $i$ of $H$ is the code $\mathbf{h}_i$ and
    $\|A\|_F = \sqrt{\sum_{i,j} a_{ij}^2}$ is the square root of the sum of
    squared entries. (We ignore the biases here.)
  \end{block}

  \begin{block}{Dimensions}
    With $m$ points in $\mathbb{R}^n$ and code size $k$:
    $X \in \mathbb{R}^{m\times n}$,\;
    $H \in \mathbb{R}^{m\times k}$ (encoder output, $H = XW^{\top}$ for a
    linear encoder with $W\in\mathbb{R}^{k\times n}$),\;
    $W^{*} \in \mathbb{R}^{k\times n}$ (decoder).
    We seek the best rank-$k$ factorisation $HW^{*}$ of the data $X$.
  \end{block}

  \begin{alertblock}{Key Recognition}
    Writing $\hat{X} = HW^{*}$, the problem is
    $\min \|X - \hat{X}\|_F^2$ over rank-$k$ matrices $\hat{X}$ ---
    \emph{exactly} the best low-rank approximation problem solved by SVD
    (Eckart--Young).
  \end{alertblock}
\end{frame}

%------------------------------------------------------------------
\begin{frame}[shrink=10]{Solving with SVD}
  \begin{block}{Optimal Low-Rank Solution}
    Let $X = U\Sigma V^{\top}$ be the SVD of the centred data. By the
    Eckart--Young theorem the optimal rank-$k$ reconstruction is
    \[
      \hat{X} = U_{\cdot,\le k}\,\Sigma_{k\times k}\,V_{\cdot,\le k}^{\top}
    \]
    where $U_{\cdot,\le k}$ keeps the first $k$ columns, etc.
  \end{block}

  \begin{block}{Matching Variables}
    Comparing $\hat{X}=HW^{*}$ with $U_{\cdot,\le k}\Sigma_{k\times k}V_{\cdot,\le k}^{\top}$,
    one valid identification is
    \[
      H = U_{\cdot,\le k}\,\Sigma_{k\times k},
      \qquad
      W^{*} = V_{\cdot,\le k}^{\top}.
    \]
    We now show the corresponding \textbf{encoder} is linear and recover
    its weight matrix $W$.
  \end{block}

  \begin{block}{Recovering the Encoder Weights $W$ (so that $H = XW^{\top}$, here written $H = MX$ logic)}
    We use $X = U\Sigma V^{\top}$, $V^{\top}V = I$, $U^{\top}U = I$, the
    pseudo-inverse identity $(XX^{\top})(XX^{\top})^{-1} = I$, and the
    transpose rule $(ABC)^{\top} = C^{\top}B^{\top}A^{\top}$.
  \end{block}
\end{frame}

%------------------------------------------------------------------
\begin{frame}[shrink=10]{Deriving $H = X\,V_{\cdot,\le k}$ (the Encoder is Linear)}
  \begin{block}{Express $H$ Using Only $X$}
    Start from $H = U_{\cdot,\le k}\Sigma_{k\times k}$ but we do not want
    $U$ explicitly --- we want $H$ as a linear map of $X$. Substitute
    $U_{\cdot,\le k}\Sigma_{k\times k}
       = (XX^{\top})(XX^{\top})^{-1}U_{\cdot,\le k}\Sigma_{k\times k}$
    and $X = U\Sigma V^{\top}$:
    \begin{align*}
      H &= (XX^{\top})(XX^{\top})^{-1}\,U_{\cdot,\le k}\Sigma_{k\times k}\\
        &= (XV\Sigma^{\top}U^{\top})\bigl(U\Sigma V^{\top}V\Sigma^{\top}U^{\top}\bigr)^{-1}
            U_{\cdot,\le k}\Sigma_{k\times k}\\
        &= XV\Sigma^{\top}U^{\top}
            \bigl(U\Sigma\Sigma^{\top}U^{\top}\bigr)^{-1}
            U_{\cdot,\le k}\Sigma_{k\times k}\\
        &= XV\Sigma^{\top}\,U^{\top}U
            (\Sigma\Sigma^{\top})^{-1}U^{\top}
            U_{\cdot,\le k}\Sigma_{k\times k}\\
        &= XV\,\Sigma^{\top}(\Sigma\Sigma^{\top})^{-1}\,
            \underbrace{U^{\top}U_{\cdot,\le k}}_{I_{\cdot,\le k}}\Sigma_{k\times k}\\
        &= XV\,\Sigma^{-1}I_{\cdot,\le k}\,\Sigma_{k\times k}
         = XV\,I_{\cdot,\le k}
         = X\,V_{\cdot,\le k}.
    \end{align*}
  \end{block}

  \begin{exampleblock}{Conclusion}
    $H = X\,V_{\cdot,\le k}$ is a \textbf{linear transformation} of $X$.
    Therefore the optimal encoder is a \emph{linear} encoder with weights
    $W = V_{\cdot,\le k}$.
  \end{exampleblock}
\end{frame}

%------------------------------------------------------------------
\begin{frame}[shrink=10]{The Encoder Weights Equal the PCA Projection Matrix}
  \begin{block}{Chain of Identities}
    \begin{itemize}\setlength\itemsep{4pt}
      \item We have the encoder $W = V_{\cdot,\le k}$.
      \item From SVD: $V$ consists of the \textbf{eigenvectors of $X^{\top}X$}.
      \item From PCA: the projection matrix $P$ consists of the
            \textbf{eigenvectors of the covariance matrix}.
      \item With the normalisation
            $\hat{x}_{ij}=\frac{1}{\sqrt{m}}\bigl(x_{ij}-\frac1m\sum_k x_{kj}\bigr)$,
            the matrix $X^{\top}X$ \emph{is} the covariance matrix.
    \end{itemize}
  \end{block}

  \begin{alertblock}{Hence Proved}
    The encoder matrix $W$ of a linear autoencoder and the projection
    matrix $P$ of PCA are computed from the eigenvectors of the same
    covariance matrix --- so they are the \textbf{same}. A linear
    undercomplete autoencoder \emph{is} PCA. \hfill $\square$
  \end{alertblock}

  \begin{block}{Remember --- the Four Conditions}
    Linear encoder, linear decoder, squared-error loss, and the
    $\frac{1}{\sqrt m}$ zero-mean normalisation. Drop any one (e.g.\ use a
    sigmoid encoder) and the equivalence no longer holds.
  \end{block}
\end{frame}

%==================================================================
\section{Regularisation in Autoencoders}
%==================================================================

\begin{frame}[shrink=10]{Why Regularisation is Needed}
  \begin{block}{Overfitting and Capacity}
    More learnable parameters $\Rightarrow$ more overfitting
    $\Rightarrow$ less generalisation. An overcomplete autoencoder
    ($\dim(\mathbf{h}) \geq \dim(\mathbf{x})$) has many parameters and can
    simply learn to \textbf{copy} $\mathbf{x}_i \to \mathbf{h} \to \hat{\mathbf{x}}_i$.
  \end{block}

  \begin{block}{Even Undercomplete Autoencoders Can Overfit}
    Consider a $100\times100$ image: $n \approx 10^4$ inputs with a code
    $k = 1000$. Although $k < n$, the model \emph{still} has a very large
    number of parameters. High-dimensional inputs with much redundancy mean
    that even shrinking the input still leaves a large parameter count, so
    poor regularisation can occur even in the undercomplete case. It is a
    \emph{more serious} problem for overcomplete autoencoders.
  \end{block}

  \begin{exampleblock}{The Fix}
    To avoid poor generalisation we introduce \textbf{regularisation}. The
    simplest solution is to add an $L_2$ penalty to the objective:
    \[
      \min_{\theta,W,W^{*},\mathbf{b},\mathbf{c}}
        \frac{1}{m}\sum_{i=1}^m\sum_{j=1}^n
        (\hat{x}_{ij}-x_{ij})^2
        \;+\; \boxed{\lambda\|\theta\|^2}
    \]
    The penalty stops weights from growing; it does not give the model
    enough freedom to overfit / memorise training data, so we hope it
    generalises well on test data.
  \end{exampleblock}
\end{frame}

%------------------------------------------------------------------
\begin{frame}[shrink=10]{Weight Tying}
  \begin{block}{Tie Encoder and Decoder Weights}
    A second trick: force the decoder weights to be the transpose of the
    encoder weights,
    \[
      \boxed{\;W^{*} = W^{\top}\;}
      \qquad (W \in \mathbb{R}^{k\times d},\; W^{*}\in\mathbb{R}^{d\times k}).
    \]
    This is enforced as a hard constraint: the same weight matrix is
    learned in such a way that it works at \emph{both} the encoder and the
    decoder.
  \end{block}

  \begin{block}{Why It Regularises}
    Tying \textbf{halves} the number of free parameters
    ($2dk \to dk$). Fewer parameters $\Rightarrow$ may generalise better.
    Geometrically, the encoder projects onto feature directions and the
    decoder reconstructs from those \emph{same} directions; this reduces
    freedom and acts as regularisation. It is commonly used for
    autoencoders.
  \end{block}

  \begin{exampleblock}{Implementation Note}
    $L_2$ regularisation is trivial to implement: it just adds a term
    $\lambda w$ (or $\alpha w$) to the gradient $\partial L/\partial w$,
    and similarly for the other parameters.
  \end{exampleblock}
\end{frame}

%==================================================================
\section{Denoising Autoencoders}
%==================================================================

\begin{frame}[shrink=10]{The Denoising Idea}
  \begin{block}{Corrupt the Input}
    A denoising autoencoder \textbf{corrupts} the input data using a
    probabilistic process $P(\tilde{x}_{ij}\mid x_{ij})$ \emph{before}
    feeding it to the network. A simple corruption used in practice:
    \[
      P(\tilde{x}_{ij}=0\mid x_{ij}) = q,
      \qquad
      P(\tilde{x}_{ij}=x_{ij}\mid x_{ij}) = 1-q.
    \]
    With probability $q$ each input is flipped to $0$; with probability
    $1-q$ it is retained as is.
  \end{block}

  \begin{block}{The Objective Stays the Same}
    Crucially, the objective is still to reconstruct the \emph{original,
    uncorrupted} $\mathbf{x}_i$:
    \[
      \arg\min_{\theta}\;
        \frac{1}{m}\sum_{i=1}^m\sum_{j=1}^n
        \bigl(\hat{x}_{ij} - x_{ij}\bigr)^2
    \]
    The network receives $\tilde{\mathbf{x}}_i$ but is graded against
    $\mathbf{x}_i$.
  \end{block}

  \begin{alertblock}{Effect}
    This is a form of regularisation that makes the representation
    \textbf{more robust}.
  \end{alertblock}
\end{frame}

%------------------------------------------------------------------
\begin{frame}[shrink=10]{Why Denoising Helps}
  \begin{block}{The Network Can No Longer Just Copy}
    With corruption, copying $\tilde{\mathbf{x}}_i \to \mathbf{h} \to
    \hat{\mathbf{x}}_i$ no longer minimises the objective: a copied
    corrupted pixel will not match the clean target, incurring error.
  \end{block}

  \begin{block}{It Must Learn Structure}
    Instead the model must learn to reconstruct a corrupted $x_{ij}$
    \emph{correctly} by relying on its \textbf{interaction with other
    elements} of $\mathbf{x}_i$. For example, a missing pixel must be
    inferred from neighbouring pixels. It therefore has to capture the
    true characteristics of the data, not memorise it.
  \end{block}

  \begin{exampleblock}{Train vs.\ Test}
    Because copying is penalised, the model can no longer get away by
    memorising the data --- yet at test time (clean inputs) it still works,
    because it learned genuine structure during training. Denoising is thus
    a kind of regularisation.
  \end{exampleblock}
\end{frame}

%------------------------------------------------------------------
\begin{frame}[shrink=10]{Application: Handwritten Digit Recognition}
  \begin{block}{Setup}
    MNIST-style digits: each image is $28\times28 = 784$ pixels, so
    $\mathbf{x}_i \in \mathbb{R}^{784}$, and labels are $0,\ldots,9$.
  \end{block}

  \begin{block}{Basic Approach vs.\ Autoencoder Approach}
    \begin{itemize}\setlength\itemsep{4pt}
      \item \textbf{Basic:} feed the $784$ raw pixels directly to a
            classifier. Requires manual feature engineering and is hard.
      \item \textbf{Autoencoder approach:} first learn the important
            characteristics of the data with an autoencoder, obtaining a
            \emph{better, smaller} representation $\mathbf{h}\in\mathbb{R}^{d}$
            ($d < 784$). No feature engineering needed.
    \end{itemize}
  \end{block}

  \begin{exampleblock}{Pipeline}
    Train the autoencoder $\to$ reuse the encoder neural network to map
    each image to its code $\mathbf{h}$ $\to$ feed the better
    representation to a multi-class SVM. Using a good learned
    representation makes the downstream classification far easier than
    feeding raw pixels. (Downside: less explainability.)
  \end{exampleblock}
\end{frame}

%------------------------------------------------------------------
\begin{frame}[shrink=10]{Visualising an Autoencoder: Neurons as Filters}
  \begin{block}{A Neuron is a Filter}
    Each hidden neuron can be viewed as a \textbf{filter} that fires (gets
    maximally activated) for a certain configuration of the input. For the
    first hidden neuron (ignoring bias),
    \[
      h_1 = \sigma\bigl(\mathbf{w}_1^{\top}\mathbf{x}_i\bigr),
    \]
    where $\mathbf{w}_1$ is the trained weight vector connecting to that
    neuron.
  \end{block}

  \begin{block}{Which Input Maximises a Neuron?}
    Training is done, so $\mathbf{w}_1$ is fixed. We ask: \emph{what input
    $\mathbf{x}_i$ maximally fires this neuron?} Assume inputs are
    normalised, $\|\mathbf{x}_i\|=1$:
    \[
      \max_{\mathbf{x}_i}\; \mathbf{w}_1^{\top}\mathbf{x}_i
      \quad\text{s.t.}\quad
      \|\mathbf{x}_i\|^2 = \mathbf{x}_i^{\top}\mathbf{x}_i = 1.
    \]
    A dot product $\mathbf{w}_1^{\top}\mathbf{x}_i =
    \|\mathbf{w}_1\|\|\mathbf{x}_i\|\cos\theta$ is maximised when
    $\cos\theta = 1$, i.e.\ both vectors point in the \textbf{same
    direction}.
  \end{block}
\end{frame}

%------------------------------------------------------------------
\begin{frame}[shrink=10]{The Maximally-Activating Input}
  \begin{block}{Solving the Constrained Maximisation}
    With the constraint $\|\mathbf{x}_i\|=1$, Cauchy--Schwarz gives
    $\mathbf{w}_1^{\top}\mathbf{x}_i \leq \|\mathbf{w}_1\|$, with equality
    when $\mathbf{x}_i \propto \mathbf{w}_1$. Solving for $\mathbf{x}_i =
    c\,\mathbf{w}_1$ under $\|\mathbf{x}_i\|=1$ gives $c = 1/\|\mathbf{w}_1\|$,
    so
    \[
      \boxed{\;
        \mathbf{x}_i = \frac{\mathbf{w}_1}{\sqrt{\mathbf{w}_1^{\top}\mathbf{w}_1}}
      \;}
    \]
    More generally, the inputs
    $\mathbf{x}_i = \dfrac{\mathbf{w}_l}{\sqrt{\mathbf{w}_l^{\top}\mathbf{w}_l}}$
    cause hidden neurons $l=1,\ldots,k$ to maximally fire.
  \end{block}

  \begin{exampleblock}{Reading the Filters}
    Plotting these maximally-activating images (reshaped to $28\times28$)
    reveals the patterns each neuron detects. For a vanilla autoencoder
    these resemble pen-strokes / parts of digits; comparing them across
    different autoencoders shows what each model has learned.
  \end{exampleblock}
\end{frame}

%------------------------------------------------------------------
\begin{frame}[shrink=10]{Denoising Filters: Pen-Stroke Detectors}
  \begin{block}{Empirical Observation}
    The hidden neurons of a denoising autoencoder act like
    \textbf{pen-stroke detectors}. For example, a highlighted neuron may
    respond to a stroke in the back region that one would expect in a
    `$0$', `$2$', `$3$', `$8$' or `$9$'.
  \end{block}

  \begin{block}{Effect of Noise Level}
    As the corruption noise $q$ \emph{increases}, the learned filters
    become \textbf{wider}. The reason: the neuron must rely on a larger set
    of adjacent pixels to feel confident about a stroke, since any single
    pixel may have been corrupted. A vanilla (no-noise) autoencoder does
    not learn many such interpretable patterns and relies more heavily on a
    few weights.
  \end{block}

  \begin{exampleblock}{Gaussian Corruption}
    Besides flipping pixels to zero, another corruption process adds
    Gaussian noise to the input:
    \[
      \tilde{x}_{ij} = x_{ij} + \mathcal{N}(0,1).
    \]
    Empirically, with Gaussian noise the AE hidden neurons behave like
    edge detectors, whereas a weight-decay-only filter does not capture
    anything special like edges.
  \end{exampleblock}
\end{frame}

%==================================================================
\section{Sparse Autoencoders}
%==================================================================

\begin{frame}[shrink=10]{Sparsity: Keep Most Neurons Inactive}
  \begin{block}{Activation of a Sigmoid Neuron}
    A hidden neuron with sigmoid activation outputs a value in $(0,1)$. We
    say it is \textbf{activated} when its output is close to $1$ and
    \textbf{inactive} when close to $0$. A \emph{sparse} autoencoder tries
    to ensure each neuron is inactive most of the time.
  \end{block}

  \begin{block}{Average Activation}
    The average activation of hidden unit $l$ over the dataset is
    \[
      \hat{\rho}_l = \frac{1}{m}\sum_{i=1}^{m} h(\mathbf{x}_i)_l .
    \]
    If unit $l$ is sparse (mostly inactive) then $\hat{\rho}_l \to 0$.
  \end{block}

  \begin{alertblock}{The Constraint}
    A sparse autoencoder uses a small \textbf{sparsity parameter} $\rho$
    (typically close to zero, e.g.\ $0.005$) and enforces
    $\boxed{\hat{\rho}_l = \rho}$ for every hidden unit $l$.
    Intuition: if a neuron fires for \emph{every} input it carries no
    discriminative power; forcing it to fire only for some inputs makes it
    capture a specific pattern. Any kind of regularisation restricts the
    freedom of parameters --- overfitting happens because of too much
    freedom.
  \end{alertblock}
\end{frame}

%------------------------------------------------------------------
\begin{frame}[shrink=10]{The Sparsity (KL-Divergence) Penalty}
  \begin{block}{Adding a Term to the Objective}
    To enforce $\hat{\rho}_l = \rho$ we add the following regularisation
    term (a sum of KL divergences) to the loss:
    \[
      \Omega(\theta)
        = \sum_{l=1}^{k}
          \rho\log\frac{\rho}{\hat{\rho}_l}
          + (1-\rho)\log\frac{1-\rho}{1-\hat{\rho}_l}
    \]
    The full regularised objective is
    \[
      \boxed{\;\widetilde{\mathcal L}(\theta) = \mathcal L(\theta) + \Omega(\theta)\;}
    \]
    where $\mathcal L(\theta)$ is the usual squared-error / cross-entropy
    reconstruction loss.
  \end{block}

  \begin{block}{Minimum of the Penalty}
    $\Omega$ reaches its minimum value (of $0$) exactly when
    $\boxed{\hat{\rho}_l = \rho}$ for all $l$ --- it is the KL divergence
    between two Bernoulli distributions, which is $0$ iff they coincide,
    and grows sharply as $\hat{\rho}_l$ departs from $\rho$.
  \end{block}

  \begin{exampleblock}{KL Divergence Intuition}
    KL divergence measures how much information is lost when one
    probability ($\hat{\rho}_l$, the neuron's empirical firing rate) is
    used to approximate another ($\rho$, the desired tiny firing rate).
  \end{exampleblock}
\end{frame}

%------------------------------------------------------------------
\begin{frame}[shrink=10]{Back-Propagating the Sparsity Penalty}
  \begin{block}{Derivative w.r.t.\ the Average Activation}
    Writing $-\Omega$ in expanded form and differentiating w.r.t.\
    $\hat{\rho}_l$:
    \[
      \frac{\partial\,\Omega(\theta)}{\partial \hat{\rho}_l}
        = -\frac{\rho}{\hat{\rho}_l}
          + \frac{1-\rho}{1-\hat{\rho}_l}.
    \]
  \end{block}

  \begin{block}{Chain Rule to the Weights}
    For each hidden neuron $l$ we need
    $\dfrac{\partial\Omega}{\partial W}$ via
    $\hat{\rho}_l \to h_i \to W$:
    \[
      \frac{\partial\Omega(\theta)}{\partial W}
        = \frac{\partial\Omega(\theta)}{\partial\hat{\rho}_l}\,
          \frac{\partial\hat{\rho}_l}{\partial W}.
    \]
    For a single weight $W_{jl}$, with
    $\hat{\rho}_l = \frac1m\sum_{i} g(\mathbf{w}_l^{\top}\mathbf{x}_i + b_l)$,
    \[
      \frac{\partial\hat{\rho}_l}{\partial W_{jl}}
        = \frac{1}{m}\sum_{i=1}^{m}
          g'\!\bigl(\mathbf{w}_l^{\top}\mathbf{x}_i + b_l\bigr)\,x_{ij},
      \quad\Rightarrow\quad
      \frac{\partial\hat{\rho}_l}{\partial \mathbf{w}_l}
        = \mathbf{x}_i\bigl(g'(\mathbf{w}_l^{\top}\mathbf{x}_i+b_l)\bigr)^{\!\top}.
    \]
  \end{block}

  \begin{alertblock}{Final Gradient}
    \[
      \frac{\partial\,\widetilde{\mathcal L}(\theta)}{\partial W}
        = \frac{\partial\,\mathcal L(\theta)}{\partial W}
          + \frac{\partial\,\Omega(\theta)}{\partial W}.
    \]
    Sparsity prevents the network from memorising: even with many hidden
    neurons, only a few are allowed to fire for any given input, so the
    model is forced to learn meaningful patterns.
  \end{alertblock}
\end{frame}

%==================================================================
\section{Contractive Autoencoders}
%==================================================================

\begin{frame}[shrink=10]{The Contractive Penalty}
  \begin{block}{Goal}
    A contractive autoencoder also tries to prevent an overcomplete
    autoencoder from learning the identity function. It does so by adding a
    \textbf{Jacobian} penalty to the loss:
    \[
      \Omega = \bigl\| J_{\mathbf{x}}(\mathbf{h}) \bigr\|_F^{2}
      \qquad (\text{Frobenius norm of the Jacobian}).
    \]
  \end{block}

  \begin{block}{What the Jacobian Looks Like}
    If the input has $n$ dimensions and the hidden layer has $k$
    dimensions, $J_{\mathbf{x}}(\mathbf{h}) \in \mathbb{R}^{k\times n}$:
    \[
      J_{\mathbf{x}}(\mathbf{h}) =
      \begin{bmatrix}
        \dfrac{\partial h_1}{\partial x_1} & \cdots & \dfrac{\partial h_1}{\partial x_n}\\[6pt]
        \dfrac{\partial h_2}{\partial x_1} & \cdots & \dfrac{\partial h_2}{\partial x_n}\\[6pt]
        \vdots & & \vdots\\[2pt]
        \dfrac{\partial h_k}{\partial x_1} & \cdots & \dfrac{\partial h_k}{\partial x_n}
      \end{bmatrix}
    \]
    The $(l,j)$ entry captures the variation in the output of the $l$-th
    neuron with a small variation in the $j$-th input.
  \end{block}

  \begin{block}{The Penalty in Sum Form}
    \[
      \bigl\| J_{\mathbf{x}}(\mathbf{h}) \bigr\|_F^{2}
        = \sum_{j=1}^{n}\sum_{l=1}^{k}
          \Bigl(\frac{\partial h_l}{\partial x_j}\Bigr)^{\!2},
      \qquad
      \min_{\theta}\; \mathcal L(\theta) + \Omega(\theta).
    \]
  \end{block}
\end{frame}

%------------------------------------------------------------------
\begin{frame}[shrink=10]{Intuition: Two Contradicting Objectives}
  \begin{block}{What $\partial h_l/\partial x_j = 0$ Means}
    A small Jacobian entry $\partial h_l/\partial x_j \approx 0$ means the
    neuron is \textbf{not sensitive} to variations in input $x_j$. The
    contractive penalty pushes the hidden representation to be insensitive
    to input changes --- it wants $\mathbf{h}$ to stay put if we perturb
    $\mathbf{x}$.
  \end{block}

  \begin{alertblock}{The Apparent Contradiction}
    But the reconstruction loss $\mathcal L(\theta)$ \emph{requires}
    $\mathbf{h}$ to capture variation in the input (otherwise it cannot
    reconstruct). So we have \textbf{two contradicting objectives}:
    \begin{itemize}\setlength\itemsep{2pt}
      \item $\mathcal L(\theta)$: \emph{capture} important variation in the data,
      \item $\Omega(\theta)$: \emph{do not} capture variation in the data.
    \end{itemize}
  \end{alertblock}

  \begin{block}{Resolution: Capture Only Important Variation}
    By pitting these two objectives against each other, we ensure
    $\mathbf{h}$ is sensitive to \emph{only the very important variations}
    observed in the training data, and insensitive to unimportant
    (noisy) directions.
    \[
      \boxed{\text{Trade-off: capture only the very important variation in the data.}}
    \]
  \end{block}
\end{frame}

%------------------------------------------------------------------
\begin{frame}[shrink=10]{Contractive AE: The PCA Geometry}
  \begin{block}{Sensitive Along High-Variance Directions Only}
    Consider the variation in the data along two directions $\mathbf{u}_1$
    (high variance) and $\mathbf{u}_2$ (low variance, essentially noise):
    \begin{itemize}\setlength\itemsep{3pt}
      \item It makes sense to let a neuron be \textbf{sensitive} to
            variation along $\mathbf{u}_1$ (important for reconstruction).
      \item It makes sense to \textbf{inhibit} sensitivity to variation
            along $\mathbf{u}_2$ (small, noisy, unimportant).
    \end{itemize}
    By doing so we balance the contradicting goals of good reconstruction
    and low sensitivity.
  \end{block}

  \begin{exampleblock}{Summary of the Regularised Autoencoders}
    \begin{itemize}\setlength\itemsep{2pt}
      \item \textbf{PCA $\equiv$ linear AE:} $\min\|X-HW^{*}\|_F^2$,
            $H = X V_{\cdot,\le k}$, related to the SVD $X = U\Sigma V^{\top}$.
      \item \textbf{$L_2$ / weight decay:} $\Omega(\theta) = \lambda\|\theta\|^2$,
            tie weights $W^{*}=W^{\top}$.
      \item \textbf{Sparse:} $\Omega(\theta) = \sum_l \rho\log\frac{\rho}{\hat\rho_l}
            + (1-\rho)\log\frac{1-\rho}{1-\hat\rho_l}$.
      \item \textbf{Denoising:} corrupt with $P(\tilde{x}_{ij}\mid x_{ij})$.
      \item \textbf{Contractive:} $\Omega(\theta)
            = \sum_{j,l}\bigl(\partial h_l/\partial x_j\bigr)^2$.
    \end{itemize}
  \end{exampleblock}
\end{frame}

%==================================================================
\section{The Bias--Variance Trade-off}
%==================================================================

\begin{frame}[shrink=10]{Setup: Fitting Two Models}
  \begin{block}{The Experiment}
    We fit a curve to a set of points sampled from a true (unknown)
    sinusoidal function. We compare two models:
    \[
      \text{Simple: } y = \hat f(x) = w_1 x + w_0 \;(\text{degree }1),
      \qquad
      \text{Complex: } y = \hat f(x) = \sum_{i=1}^{25} w_i x^i + w_0 \;(\text{degree }25).
    \]
    In both cases we make an assumption about the form, with no idea of the
    true relation $f(x)$.
  \end{block}

  \begin{block}{Repeated Sampling}
    We sample (say) $25$ points from the training data and train a simple
    and a complex model. We repeat $k$ times --- each model is trained on a
    \emph{different} sample of training data.
  \end{block}

  \begin{exampleblock}{Observations}
    \begin{itemize}\setlength\itemsep{3pt}
      \item \textbf{Simple models} trained on different samples do not
            differ much from each other (low spread), but all lie far from
            the true sinusoid $\Rightarrow$ \textbf{underfitting}.
      \item \textbf{Complex models} trained on different samples are very
            different from each other (high spread) $\Rightarrow$
            \textbf{overfitting}; different data $\Rightarrow$ very
            different output.
    \end{itemize}
  \end{exampleblock}
\end{frame}

%------------------------------------------------------------------
\begin{frame}[shrink=10]{Defining Bias and Variance}
  \begin{block}{Bias}
    Let $f(x)$ be the true model and $\hat f(x)$ our estimate. The
    \textbf{bias} measures how far the \emph{average} model is from the
    truth:
    \[
      \mathrm{Bias}\bigl(\hat f(x)\bigr)
        = \mathbb{E}\bigl[\hat f(x)\bigr] - f(x).
    \]
    The simple model's average (blue line) is far from the true sinusoid
    $\Rightarrow$ \textbf{high bias}. The complex model's average is close
    $\Rightarrow$ \textbf{low bias}.
  \end{block}

  \begin{block}{Variance}
    The \textbf{variance} measures how much the estimate spreads across
    different training samples:
    \[
      \mathrm{Var}\bigl(\hat f(x)\bigr)
        = \mathbb{E}\Bigl[\bigl(\hat f(x) - \mathbb{E}[\hat f(x)]\bigr)^2\Bigr].
    \]
    Roughly, it tells us how much the different $\hat f(x)$'s (trained on
    different samples) differ from each other. Simple model: \textbf{low
    variance}; complex model: \textbf{high variance}.
  \end{block}

  \begin{alertblock}{Informal Summary}
    Simple model: \textbf{high bias, low variance}.\quad
    Complex model: \textbf{low bias, high variance}.\quad
    There is always a \textbf{trade-off}; both bias and variance contribute
    to the mean-squared error (MSE).
  \end{alertblock}
\end{frame}

%------------------------------------------------------------------
\begin{frame}[shrink=10]{The Bias--Variance Decomposition}
  \begin{block}{Statement}
    For a new point with $y = f(x)+\varepsilon$, $\varepsilon\sim\mathcal N(0,\sigma^2)$,
    the expected squared prediction error decomposes as
    \[
      \boxed{\;
        \mathbb{E}\bigl[(y - \hat f(x))^2\bigr]
          = \underbrace{\mathrm{Bias}^2}_{}
            + \underbrace{\mathrm{Variance}}_{}
            + \underbrace{\sigma^2}_{\text{irreducible noise}}
      \;}
    \]
  \end{block}

  \begin{block}{Proof Sketch}
    Write $y - \hat f = \bigl(f - \mathbb{E}[\hat f]\bigr)
    + \bigl(\mathbb{E}[\hat f] - \hat f\bigr) + \varepsilon$.
    Expanding the square and taking expectations, the cross terms vanish
    because (i) $\mathbb{E}[\varepsilon]=0$ and $\varepsilon$ is independent
    of $\hat f$, and (ii) $\mathbb{E}\bigl[\mathbb{E}[\hat f]-\hat f\bigr]=0$.
    Using
    \[
      A = f - \mathbb{E}[\hat f],\quad
      B = \varepsilon,\quad
      C = \mathbb{E}[\hat f] - \hat f,
    \]
    only $\mathbb{E}[A^2]=\mathrm{Bias}^2$, $\mathbb{E}[B^2]=\sigma^2$,
    $\mathbb{E}[C^2]=\mathrm{Variance}$ survive.
  \end{block}

  \begin{exampleblock}{Conclusion}
    \[
      \text{Expected error} = \mathrm{Bias}^2 + \mathrm{Variance}
        + \text{Irreducible Noise}.
    \]
    Reducing bias (more complexity) increases variance and vice versa ---
    we must balance the two.
  \end{exampleblock}
\end{frame}

%==================================================================
\section{Training Error vs Test Error}
%==================================================================

\begin{frame}[shrink=10]{The Two Errors}
  \begin{block}{Train Error and Test Error}
    The parameters of $\hat f(x)$ are trained on a training set
    $\{(x_i,y_i)\}_{i=1}^{n}$. At test time we evaluate on a
    \emph{validation/unseen} set not used for training. This gives two
    quantities:
    \[
      \text{train}_{\text{err}} = \frac{1}{n}\sum_{i=1}^{n}\bigl(y_i - \hat f(x_i)\bigr)^2,
      \quad
      \text{test}_{\text{err}} = \frac{1}{m}\sum_{i=n+1}^{n+m}\bigl(y_i - \hat f(x_i)\bigr)^2.
    \]
    As model complexity rises, train error becomes overly optimistic and
    gives a wrong picture of how close $\hat f$ is to $f$. The validation
    error gives the real picture.
  \end{block}

  \begin{block}{True Error}
    For an unseen point we want the \textbf{true error}
    $\mathbb{E}\bigl[(\hat f(x_i) - f(x_i))^2\bigr]$. We cannot estimate it
    directly because we do not know $f$. We estimate it empirically using
    observations $y_i$ and predictions $\hat y_i$.
  \end{block}

  \begin{alertblock}{Model Used}
    $D = \{(x_i,y_i)\}_{i=1}^{m+n}$ with $y_i = f(x_i) + \varepsilon_i$,
    $\varepsilon \sim \mathcal N(0,\sigma^2)$, and $\hat f$ trained on
    $T \subset D$ (training data), giving predictions $\hat y_i = \hat f(x_i)$.
  \end{alertblock}
\end{frame}

%------------------------------------------------------------------
\begin{frame}[shrink=10]{Relating Empirical Error to True Error}
  \begin{block}{Empirical Expectation}
    The empirical squared error over observations is
    \[
      \mathbb{E}\bigl[(\hat y_i - y_i)^2\bigr]
        = \frac{1}{m}\sum_{i=1}^{m}(\hat y_i - y_i)^2.
    \]
    (Analogy: to estimate expected goals $\mathbb{E}[\text{goals}]$ over
    $k$ matches we use $\frac1k\sum_i g_i$.)
  \end{block}

  \begin{block}{Substituting $y_i = f(x_i)+\varepsilon_i$}
    Expanding $\mathbb{E}[(\hat y_i - y_i)^2]$ with $\hat y_i = \hat f(x_i)$:
    \begin{align*}
      \mathbb{E}\bigl[(\hat f(x_i) - f(x_i))^2\bigr]
        &= \mathbb{E}\bigl[(\hat y_i - y_i)^2\bigr]
           - \mathbb{E}[\varepsilon_i^2]
           + 2\,\mathbb{E}\bigl[\varepsilon_i\,(\hat f(x_i) - f(x_i))\bigr].
    \end{align*}
    We can empirically evaluate the first RHS term using observations. The
    last term is a \textbf{covariance} between $\varepsilon_i$ and
    $\hat f(x_i)-f(x_i)$ (since $\mathbb{E}[\varepsilon_i]=0$).
  \end{block}

  \begin{exampleblock}{Covariance Reminder}
    $\mathrm{Cov}(x,y) = \mathbb{E}[(x-\mu_x)(y-\mu_y)] = \mathbb{E}[xy] -
    \mu_x\mu_y$; with $\mu_{\varepsilon}=0$ the cross term is exactly
    $\mathbb{E}[\varepsilon_i(\hat f(x_i)-f(x_i))]$.
  \end{exampleblock}
\end{frame}

%------------------------------------------------------------------
\begin{frame}[shrink=10]{Case 1 --- Using Test Observations}
  \begin{block}{The Key Independence}
    With \textbf{test} data, none of the test observations participated in
    estimating $\hat f$ (its parameters were learned only on training
    data). Therefore the test noise $\varepsilon_i$ is \textbf{independent}
    of $\hat f(x_i) - f(x_i)$:
    \[
      \varepsilon_i \perp \bigl(\hat f(x_i) - f(x_i)\bigr).
    \]
  \end{block}

  \begin{block}{The Cross Term Vanishes}
    For independent variables $\mathbb{E}[XY] = \mathbb{E}[X]\mathbb{E}[Y]$,
    so
    \[
      \mathbb{E}\bigl[\varepsilon_i\,(\hat f(x_i)-f(x_i))\bigr]
        = \underbrace{\mathbb{E}[\varepsilon_i]}_{=0}\,
          \mathbb{E}\bigl[\hat f(x_i)-f(x_i)\bigr]
        = 0.
    \]
  \end{block}

  \begin{alertblock}{Result}
    \[
      \boxed{\;\text{true error} = \text{empirical test error} + \text{small constant}\;}
    \]
    i.e.\ $\mathbb{E}[(\hat f(x_i)-f(x_i))^2]
      = \frac{1}{m}\sum_{i=n+1}^{n+m}(\hat y_i - y_i)^2 - \frac1m\sum\varepsilon_i^2$.
    The test error is close to the true error, differing only by a small
    constant. That is why we use a validation set.
  \end{alertblock}
\end{frame}

%------------------------------------------------------------------
\begin{frame}[shrink=10]{Case 2 --- Using Training Observations (Optimism)}
  \begin{block}{No Independence Here}
    With \textbf{training} data, $\varepsilon_i$ \emph{was} used to estimate
    the parameters of $\hat f(x)$. Hence $\varepsilon_i$ is correlated with
    $\hat f(x_i)$, and
    \[
      \mathbb{E}\bigl[\varepsilon_i\,(\hat f(x_i)-f(x_i))\bigr]
        \neq \mathbb{E}[\varepsilon_i]\,\mathbb{E}\bigl[\hat f(x_i)-f(x_i)\bigr]
        \neq 0.
    \]
  \end{block}

  \begin{block}{Consequence: Optimism}
    The covariance term is positive, so the empirical \textbf{train error
    is smaller} than the true error --- it does not give the true picture
    of the error. The training error is \textbf{overly optimistic}.
  \end{block}

  \begin{exampleblock}{Practical Rule}
    Even gradient-based optimisers (SGD, Adam) make stationary
    distribution assumptions about the noise; if those hold, the noise is
    zero-mean and the analysis carries through. Regardless,
    \[
      \boxed{\text{Always use a validation set (independent of training) to estimate error.}}
    \]
  \end{exampleblock}
\end{frame}

%==================================================================
\section{True Error and Model Complexity}
%==================================================================

\begin{frame}[shrink=10]{Quantifying the Optimism via Stein's Lemma}
  \begin{block}{Stein's Lemma Result}
    Using Stein's lemma (with some algebra) one can show the train-error
    covariance term equals
    \[
      \boxed{\;
        \frac{1}{n}\sum_{i=1}^{n}\varepsilon_i\bigl(\hat f(x_i) - f(x_i)\bigr)
          = \frac{\sigma^2}{n}\sum_{i=1}^{n}\frac{\partial \hat f(x_i)}{\partial y_i}
      \;}
    \]
    (Stirling's approximation $n! \approx \sqrt{2\pi n}\,(n/e)^n$ is used in
    the derivation of related results.)
  \end{block}

  \begin{block}{When is $\partial\hat f(x_i)/\partial y_i$ Large?}
    It is large when a \emph{small} change in the observation $y_i$ causes a
    \emph{large} change in the estimate $\hat f$. A \textbf{complex} model
    is more sensitive to changes in the observations; a \textbf{simple}
    model is less sensitive. So this derivative directly measures model
    complexity.
  \end{block}

  \begin{exampleblock}{Verification}
    Fit a simple and a complex model to data, then change one data point
    and retrain. The simple model barely changes; the complex model changes
    a lot --- confirming complex models are more sensitive.
  \end{exampleblock}
\end{frame}

%------------------------------------------------------------------
\begin{frame}[shrink=10]{The Master Equation of Regularisation}
  \begin{block}{True Error = Train Error + Complexity}
    Combining the pieces:
    \[
      \boxed{\;
        \text{true error}
          = \text{empirical train error}
            + \text{small constant}
            + \Omega(\text{model complexity})
      \;}
    \]
    where $\Omega(\theta)$ approximates
    $\dfrac{\sigma^2}{n}\sum_{i=1}^{n}\dfrac{\partial\hat f(x_i)}{\partial y_i}$;
    it is high for complex models, small for simple models.
  \end{block}

  \begin{block}{What We Should Actually Minimise}
    Instead of minimising the training error $\mathcal L_{\text{train}}(\theta)$
    alone, we should minimise
    \[
      \min_{\theta}\;
        \underbrace{\mathcal L_{\text{train}}(\theta) + \Omega(\theta)}_{\mathcal L(\theta)}.
    \]
    $\Omega(\theta)$ should ensure the model has \emph{reasonable}
    complexity, so the model generalises best on test data. This is the
    \textbf{basis for all regularisation methods}.
  \end{block}

  \begin{alertblock}{Why We Care (Universal Approximation)}
    By the Universal Approximation Theorem, deep neural networks are highly
    complex models with many parameters and non-linearities; it is easy for
    them to overfit and drive training error to zero. Hence we need some
    form of regularisation.
  \end{alertblock}
\end{frame}

%------------------------------------------------------------------
\begin{frame}[shrink=10]{The Sweet Spot}
  \begin{block}{Train Error vs Test Error vs Complexity}
    As model complexity increases:
    \begin{itemize}\setlength\itemsep{3pt}
      \item \textbf{Train error} decreases monotonically (towards $0$).
      \item \textbf{Test error} first decreases (less bias) then increases
            (more variance) --- a U-shape.
    \end{itemize}
    The minimum of the test-error curve is the \textbf{sweet spot}: the
    ideal model complexity, achieved by tuning $\Omega(\theta)$.
  \end{block}

  \begin{exampleblock}{Different Forms of Regularisation}
    All of the following are instances of adding an $\Omega(\theta)$ that
    controls complexity:
    \begin{itemize}\setlength\itemsep{1pt}
      \item $L_2$ regularisation (weight decay)
      \item Dataset augmentation
      \item Parameter sharing and tying
      \item Adding noise to the inputs
      \item Adding noise to the outputs (label smoothing)
      \item Early stopping
      \item Ensemble methods
      \item Dropout
    \end{itemize}
    We treat each in turn.
  \end{exampleblock}
\end{frame}

%==================================================================
\section{$L_2$ Regularisation}
%==================================================================

\begin{frame}[shrink=10]{$L_2$ Regularisation: Setup}
  \begin{block}{Regularised Loss}
    \[
      \widetilde{\mathcal L}(\mathbf{w})
        = \underbrace{\mathcal L(\mathbf{w})}_{\text{empirical train error}}
          + \frac{\alpha}{2}\|\mathbf{w}\|^2 .
    \]
    The penalty relates to model complexity by \emph{reducing the freedom}
    of the weights.
  \end{block}

  \begin{block}{Gradient and Update}
    \[
      \nabla\widetilde{\mathcal L}(\mathbf{w}) = \nabla\mathcal L(\mathbf{w}) + \alpha\mathbf{w},
      \qquad
      \mathbf{w}_{t+1} = \mathbf{w}_t - \eta\nabla\mathcal L(\mathbf{w}_t) - \eta\alpha\mathbf{w}_t.
    \]
    This requires only a very small modification to the code: subtract an
    extra $\eta\alpha\mathbf{w}_t$ each step (weight decay).
  \end{block}

  \begin{block}{Taylor Expansion around the Unregularised Optimum}
    Let $\mathbf{w}^{*}$ be optimal for the \emph{unregularised} loss, so
    $\nabla\mathcal L(\mathbf{w}^{*}) = 0$. Expanding to second order with
    Hessian $H$:
    \[
      \mathcal L(\mathbf{w}) \approx \mathcal L(\mathbf{w}^{*})
        + \tfrac12(\mathbf{w}-\mathbf{w}^{*})^{\top}H(\mathbf{w}-\mathbf{w}^{*}),
      \qquad
      \nabla\mathcal L(\mathbf{w}) = H(\mathbf{w}-\mathbf{w}^{*}).
    \]
  \end{block}
\end{frame}

%------------------------------------------------------------------
\begin{frame}[shrink=10]{$L_2$: Solving for the Regularised Optimum}
  \begin{block}{Set the Regularised Gradient to Zero}
    Let $\tilde{\mathbf{w}}$ be optimal for $\widetilde{\mathcal L}$, so
    $\nabla\widetilde{\mathcal L}(\tilde{\mathbf{w}}) = 0$:
    \[
      \nabla\mathcal L(\tilde{\mathbf{w}}) + \alpha\tilde{\mathbf{w}}
        = H(\tilde{\mathbf{w}}-\mathbf{w}^{*}) + \alpha\tilde{\mathbf{w}} = 0
      \;\Rightarrow\;
      (H+\alpha I)\tilde{\mathbf{w}} = H\mathbf{w}^{*}
      \;\Rightarrow\;
      \boxed{\tilde{\mathbf{w}} = (H+\alpha I)^{-1}H\mathbf{w}^{*}}
    \]
    As $\alpha\to 0$, $\tilde{\mathbf{w}}\to\mathbf{w}^{*}$ (no
    regularisation). We analyse $\alpha\neq 0$.
  \end{block}

  \begin{block}{Diagonalise the Hessian}
    Assume $H$ is symmetric PSD with $H = Q\Lambda Q^{\top}$ ($Q$
    orthogonal, $QQ^{\top}=Q^{\top}Q=I$). Then
    \begin{align*}
      \tilde{\mathbf{w}}
        &= (Q\Lambda Q^{\top} + \alpha QIQ^{\top})^{-1} Q\Lambda Q^{\top}\mathbf{w}^{*}
         = \bigl(Q(\Lambda+\alpha I)Q^{\top}\bigr)^{-1} Q\Lambda Q^{\top}\mathbf{w}^{*}\\
        &= Q(\Lambda+\alpha I)^{-1}Q^{\top}\,Q\Lambda Q^{\top}\mathbf{w}^{*}
         = Q\,\underbrace{(\Lambda+\alpha I)^{-1}\Lambda}_{D}\,Q^{\top}\mathbf{w}^{*}.
    \end{align*}
  \end{block}

  \begin{exampleblock}{Result}
    \[
      \boxed{\;
        \tilde{\mathbf{w}} = Q\,(\Lambda+\alpha I)^{-1}\Lambda\,Q^{\top}\mathbf{w}^{*}
        ,\qquad
        D = (\Lambda+\alpha I)^{-1}\Lambda \text{ is diagonal.}
      \;}
    \]
  \end{exampleblock}
\end{frame}

%------------------------------------------------------------------
\begin{frame}[shrink=10]{$L_2$: What Is Happening Geometrically}
  \begin{block}{Rotate, Scale, Rotate Back}
    The operation $\tilde{\mathbf{w}} = QDQ^{\top}\mathbf{w}^{*}$ does three
    things:
    \begin{enumerate}\setlength\itemsep{2pt}
      \item $Q^{\top}\mathbf{w}^{*}$: rotate $\mathbf{w}^{*}$ into the
            eigenbasis.
      \item $D$: scale each coordinate.
      \item $Q$: rotate back.
    \end{enumerate}
    If $\alpha=0$ then $D=I$ and $Q$ rotates $Q^{\top}\mathbf{w}^{*}$ back to
    $\mathbf{w}^{*}$ (nothing changes).
  \end{block}

  \begin{block}{The Diagonal Scaling}
    For $\alpha\neq 0$, $D = \mathrm{diag}\!\Bigl(\dfrac{d_i}{d_i+\alpha}\Bigr)$
    where $d_i$ are the eigenvalues of $H$:
    \[
      d_i \gg \alpha \;\Rightarrow\; \frac{d_i}{d_i+\alpha}\approx 1
      \quad(\text{direction kept}),\qquad
      d_i \ll \alpha \;\Rightarrow\; \frac{d_i}{d_i+\alpha}\approx 0
      \quad(\text{direction killed}).
    \]
  \end{block}

  \begin{alertblock}{Interpretation}
    The weight vector is rotated to $\tilde{\mathbf{w}}$; \emph{all}
    elements shrink, but some shrink \emph{more} than others. Important
    directions (large eigenvalue, large loss curvature) are retained;
    unimportant directions are scaled down. This is exactly
    \textbf{dimension-specific} shrinkage --- what we wanted. The effective
    number of parameters is $\sum_{i=1}^{n}\dfrac{d_i}{d_i+\alpha} \ll n$.
  \end{alertblock}
\end{frame}

%==================================================================
\section{Dataset Augmentation}
%==================================================================

\begin{frame}[shrink=10]{Dataset Augmentation}
  \begin{block}{Idea}
    Minimising the empirical train error on a small dataset leads to
    perfect classification of training points (overfitting). Augmentation
    creates \emph{more} labelled data using domain knowledge of the task.
    Given an image, we exploit the fact that certain transformations do
    \textbf{not change the label}:
    \begin{itemize}\setlength\itemsep{1pt}
      \item rotate by $20^{\circ}$, by $62^{\circ}$,
      \item shift vertically / horizontally,
      \item blur,
      \item change some pixels.
    \end{itemize}
  \end{block}

  \begin{block}{Why It Works}
    Typically \textbf{more data $\Rightarrow$ better learning}. The
    augmented data is supervised data created using knowledge of the task.
    It works very well for image classification / object recognition and
    has also been shown to work for speech.
  \end{block}

  \begin{exampleblock}{Caveat}
    For some tasks it may not be clear how to generate such label-preserving
    data (e.g.\ some NLP tasks), where naive transformations could change
    meaning. Augmentation must respect the semantics of the task.
  \end{exampleblock}
\end{frame}

%==================================================================
\section{Parameter Sharing and Tying}
%==================================================================

\begin{frame}[shrink=10]{Parameter Sharing vs.\ Parameter Tying}
  \begin{columns}[T]
    \column{0.48\textwidth}
    \begin{block}{Parameter Sharing}
      \begin{itemize}\setlength\itemsep{2pt}
        \item Used in \textbf{CNNs}.
        \item The same filter is applied at different positions of the
              image.
        \item Equivalently, the same weight matrix acts on different input
              neurons.
        \item Use the same parameter in multiple computations.
      \end{itemize}
    \end{block}

    \column{0.48\textwidth}
    \begin{block}{Parameter Tying}
      \begin{itemize}\setlength\itemsep{2pt}
        \item Typically used in \textbf{autoencoders}.
        \item Encoder and decoder weights are tied.
        \item Impose an equality constraint between parameters:
              $W^{*} = W^{\top}$.
      \end{itemize}
    \end{block}
  \end{columns}

  \medskip

  \begin{block}{Why Tying Regularises}
    Tying is a specific form of sharing that \textbf{reduces the number of
    free parameters} and acts as regularisation (it reduces model
    complexity $\Omega(\theta)$, consistent with the Stein's-lemma view).
    With $W \in \mathbb{R}^{k\times d}$ and $W^{*}\in\mathbb{R}^{d\times k}$,
    tying $W^{*}=W^{\top}$ takes the parameter count from $2dk$ down to
    $dk$ --- \textbf{half the parameters}.
  \end{block}
\end{frame}

%==================================================================
\section{Adding Noise to the Inputs}
%==================================================================

\begin{frame}[shrink=10]{Adding Gaussian Noise to Inputs $\equiv$ Weight Decay}
  \begin{block}{Setup}
    We already saw input noise in denoising autoencoders. For a simple
    linear input--output network, adding Gaussian noise to the input is
    equivalent to \textbf{weight decay} ($L_2$ regularisation). It can also
    be viewed as data augmentation.
    Let $\tilde{x}_i = x_i + \varepsilon_i$, $\varepsilon_i \sim \mathcal N(0,\sigma^2)$.
  \end{block}

  \begin{block}{Original vs.\ Noisy Output}
    \[
      \hat y = \sum_{i=1}^{n} w_i x_i \quad(\text{original}),
      \qquad
      \tilde y = \sum_{i=1}^{n} w_i \tilde{x}_i
              = \hat y + \sum_{i=1}^{n} w_i \varepsilon_i .
    \]
  \end{block}

  \begin{block}{Expected Squared Error with Noise}
    \begin{align*}
      \mathbb{E}\bigl[(\tilde y - y)^2\bigr]
        &= \mathbb{E}\Bigl[\bigl((\hat y - y) + \textstyle\sum_i w_i\varepsilon_i\bigr)^2\Bigr]\\
        &= \mathbb{E}[(\hat y - y)^2]
           + 2\,\mathbb{E}\Bigl[(\hat y - y)\textstyle\sum_i w_i\varepsilon_i\Bigr]
           + \mathbb{E}\Bigl[\bigl(\textstyle\sum_i w_i\varepsilon_i\bigr)^2\Bigr].
    \end{align*}
  \end{block}
\end{frame}

%------------------------------------------------------------------
\begin{frame}[shrink=10]{Input Noise: The Cross Term Vanishes}
  \begin{block}{Cross Term}
    Since $\mathbb{E}[\varepsilon_i]=0$ and $\varepsilon_i$ is independent
    of $\hat y - y$,
    \[
      2\,\mathbb{E}\Bigl[(\hat y - y)\sum_i w_i\varepsilon_i\Bigr]
        = 2\sum_i w_i\,\mathbb{E}[\hat y - y]\,\mathbb{E}[\varepsilon_i] = 0.
    \]
  \end{block}

  \begin{block}{Noise Term}
    Using $\mathbb{E}[\varepsilon_i\varepsilon_j]=0$ for $i\neq j$
    (independent) and $\mathbb{E}[\varepsilon_i^2]=\sigma^2$,
    \[
      \mathbb{E}\Bigl[\bigl(\textstyle\sum_i w_i\varepsilon_i\bigr)^2\Bigr]
        = \sum_{i}\sum_{j} w_i w_j\,\mathbb{E}[\varepsilon_i\varepsilon_j]
        = \sigma^2\sum_{i=1}^{n} w_i^2 .
    \]
  \end{block}

  \begin{alertblock}{Result}
    \[
      \boxed{\;
        \mathbb{E}\bigl[(\tilde y - y)^2\bigr]
          = \mathbb{E}\bigl[(\hat y - y)^2\bigr]
            + \sigma^2\sum_{i=1}^{n} w_i^2
      \;}
    \]
    The extra term $\sigma^2\sum_i w_i^2$ is exactly the $L_2$ norm penalty.
    Hence \textbf{adding noise to the input is the same as $L_2$
    regularisation}.
  \end{alertblock}
\end{frame}

%==================================================================
\section{Adding Noise to the Outputs (Label Smoothing)}
%==================================================================

\begin{frame}[shrink=10]{Hard Targets vs.\ Soft Targets}
  \begin{block}{Hard Targets}
    With one-hot labels, we minimise $\sum_{l\ge 0} p_l\log q_l$ where the
    true distribution is $p = \{0,1,0,\ldots,0\}$ (all mass on the correct
    class) and $q$ is the estimate. Trusting these hard labels completely
    leads the model to drive training error to zero --- overfitting.
  \end{block}

  \begin{block}{Soft Targets (Label Smoothing)}
    Do not trust the labels as perfectly true --- they may be noisy.
    Instead use \textbf{soft targets}: spread a small amount of probability
    mass $\varepsilon$ across the wrong classes,
    \[
      p = \Bigl\{\tfrac{\varepsilon}{k-1},\ \tfrac{\varepsilon}{k-1},\ \ldots,\ 1-\varepsilon,\ \ldots,\ \tfrac{\varepsilon}{k-1}\Bigr\},
    \]
    where $\varepsilon$ is the noise in the label (a small positive
    constant) and $k$ is the number of classes.
  \end{block}

  \begin{exampleblock}{Why It Acts Like Regularisation}
    The aim of regularisation is to not overfit the training data. Soft
    targets ensure the model treats the labels as only \emph{partly}
    reliable, giving corrupted information. The principle: \textbf{do not
    allow the true training error (computed from training data) to go to
    zero} --- do whatever you can to avoid that. This is a heuristic.
  \end{exampleblock}
\end{frame}

%==================================================================
\section{Early Stopping}
%==================================================================

\begin{frame}[shrink=10]{Early Stopping: The Algorithm}
  \begin{block}{Procedure}
    \begin{itemize}\setlength\itemsep{3pt}
      \item Track the \textbf{validation error} as training proceeds.
      \item Keep a \textbf{patience} parameter $p$.
      \item Every $p$ steps, check whether validation error has improved
            (got worse, increased, or remained the same).
      \item If at step $k$ there was \emph{no improvement} in the previous
            $p$ steps, \textbf{stop} training and return the model stored at
            step $k-p$.
    \end{itemize}
    Training error keeps decreasing, but validation error follows a U-shape
    --- we stop near its minimum.
  \end{block}

  \begin{exampleblock}{Example (patience $p=5$)}
    \centering
    \begin{tabular}{cc}
      \toprule
      \textbf{Epoch} & \textbf{Val.\ loss}\\
      \midrule
      $10$ & $0.25$\\
      $11$ & $0.23$\\
      $12$ & $0.24$\\
      $13$ & $0.25$\\
      $14$ & $0.25$\\
      $15$ & $0.24$\\
      $16$ & $0.24$\\
      \bottomrule
    \end{tabular}
  \end{exampleblock}

  \begin{block}{Properties}
    Very effective and one of the most widely used forms of regularisation.
    Can be combined with other regularisers (such as $L_2$). But \emph{how}
    does it act as a regulariser, given gradient descent only sees training
    data?
  \end{block}
\end{frame}

%------------------------------------------------------------------
\begin{frame}[shrink=10]{Early Stopping: Why It Limits Weight Growth}
  \begin{block}{Bounding the Update}
    The SGD update $\mathbf{w}_{t+1} = \mathbf{w}_t - \eta\nabla\mathbf{w}_t$
    unrolls to $\mathbf{w}_{t+1} = \mathbf{w}_0 - \eta\sum_{k=1}^{t}\nabla\mathbf{w}_k$.
    If $\tau$ is the maximum value of $|\nabla\mathbf{w}|$, then
    \[
      |\mathbf{w}_{t+1}| \;\le\; |\mathbf{w}_0| + \eta\,t\,\tau .
    \]
    So $t$ \emph{controls how far} $\mathbf{w}_t$ can travel from
    $\mathbf{w}_0$ --- limited steps means limited weight growth.
  \end{block}

  \begin{block}{Connecting to the Hessian}
    Using the Taylor expansion $\nabla\mathcal L(\mathbf{w}) = H(\mathbf{w}-\mathbf{w}^{*})$
    and starting from $\mathbf{w}_0 = 0$, the SGD update becomes
    \[
      \mathbf{w}_t = (I+\eta H)\mathbf{w}_{t-1} - \eta H\mathbf{w}^{*},
    \]
    and with $H = Q\Lambda Q^{\top}$ one can show
    \[
      \boxed{\;
        \mathbf{w}_t
          = Q\bigl(I - (I-\varepsilon\Lambda)^{t}\bigr)Q^{\top}\mathbf{w}^{*}
      \;}
      \quad(\text{rotation}\to\text{scaling}\to\text{rotation}).
    \]
  \end{block}

  \begin{alertblock}{Comparison with $L_2$}
    Recall the $L_2$ optimum
    $\tilde{\mathbf{w}} = Q\bigl(I - (\Lambda+\alpha I)^{-1}\alpha\bigr)Q^{\top}\mathbf{w}^{*}$.
    The two coincide (i.e.\ $\mathbf{w}_t = \tilde{\mathbf{w}}$) when
    $\varepsilon$, $t$ and $\alpha$ are chosen so that
    $(I-\varepsilon\Lambda)^{t} = (I+\alpha\Lambda^{-1})^{-1}$. So early
    stopping is, in effect, $L_2$ regularisation in disguise.
  \end{alertblock}
\end{frame}

%------------------------------------------------------------------
\begin{frame}[shrink=10]{Early Stopping: Direction-Wise View}
  \begin{block}{Important vs.\ Unimportant Directions}
    \begin{itemize}\setlength\itemsep{4pt}
      \item Early stopping only allows $t$ updates to the parameters.
      \item If a parameter / direction $\mathbf{w}$ is \textbf{important}
            for the loss $\mathcal L(\theta)$, then $\partial\mathcal L/\partial w$
            is large, so its updates are large and it grows quickly.
      \item If a parameter is \textbf{not important},
            $\partial\mathcal L/\partial w$ is small, its updates are small,
            and within only $t$ steps it cannot grow large.
    \end{itemize}
  \end{block}

  \begin{exampleblock}{Conclusion}
    By doing a \emph{limited} number of steps $t$, early stopping keeps
    unimportant weights small (just like $L_2$ shrinks them via
    $d_i/(d_i+\alpha)$), retaining only the significant directions. This is
    why early stopping generalises.
  \end{exampleblock}
\end{frame}

%==================================================================
\section{Ensemble Methods (Bagging)}
%==================================================================

\begin{frame}[shrink=10]{Ensembles and Bagging}
  \begin{block}{Idea}
    Combine the outputs of different models to reduce generalisation error.
    The models can respond to different classifiers, or be different
    instances of the same classifier trained with: different
    hyperparameters, different features, or different samples of training
    data.
  \end{block}

  \begin{block}{Bagging}
    Form an ensemble of different instances of the \emph{same} classifier.
    From a given dataset, construct multiple training sets
    $T_1, T_2, \ldots, T_k$ by \textbf{sampling with replacement} (each a
    subset). Train one model on each.
  \end{block}

  \begin{exampleblock}{When Would Bagging Work?}
    Each classifier makes certain errors on different points. We \emph{want
    the errors to be as independent as possible} across models --- then
    averaging cancels them out.
  \end{exampleblock}
\end{frame}

%------------------------------------------------------------------
\begin{frame}[shrink=10]{Bagging: The MSE of the Ensemble}
  \begin{block}{Setup}
    Consider $k$ models. Model $i$ makes error $\varepsilon_i$ on a test
    example, where each $\varepsilon_i$ is drawn from a zero-mean
    multivariate normal with
    \[
      \mathrm{Var} = \mathbb{E}[\varepsilon_i^2] = V,
      \qquad
      \mathrm{Cov} = \mathbb{E}[\varepsilon_i\varepsilon_j] = C \;(i\neq j).
    \]
    The average prediction error is $\frac1k\sum_i\varepsilon_i$.
  \end{block}

  \begin{block}{Expected Squared Error}
    \begin{align*}
      \text{mse}
        &= \mathbb{E}\Bigl[\bigl(\tfrac1k\textstyle\sum_i\varepsilon_i\bigr)^2\Bigr]
         = \frac{1}{k^2}\,\mathbb{E}\Bigl[\textstyle\sum_i\varepsilon_i^2
             + \sum_{i}\sum_{j\neq i}\varepsilon_i\varepsilon_j\Bigr]\\
        &= \frac{1}{k^2}\Bigl[\textstyle\sum_i\mathbb{E}[\varepsilon_i^2]
             + \sum_i\sum_{j\neq i}\mathbb{E}[\varepsilon_i\varepsilon_j]\Bigr]
         = \frac{1}{k^2}\bigl[kV + k(k-1)C\bigr]\\
        &= \boxed{\;\frac{1}{k}V + \frac{k-1}{k}C\;}
    \end{align*}
    (using $k$ terms with variance $V$ and $k(k-1)$ cross terms with
    covariance $C$).
  \end{block}
\end{frame}

%------------------------------------------------------------------
\begin{frame}[shrink=10]{Bagging: When Does It Help?}
  \begin{block}{Two Extremes}
    \begin{itemize}\setlength\itemsep{5pt}
      \item \textbf{Errors perfectly correlated} ($C = V$):
            \[
              \text{mse} = \frac{1}{k}V + \frac{k-1}{k}V = V.
            \]
            Bagging \textbf{does not help}: the ensemble MSE is as bad as a
            single model.
      \item \textbf{Errors independent / uncorrelated} ($C = 0$):
            \[
              \text{mse} = \frac{1}{k}V.
            \]
            The ensemble MSE \textbf{shrinks by a factor of $k$}.
    \end{itemize}
  \end{block}

  \begin{alertblock}{Takeaway}
    Bagging works precisely when the models' errors are independent (or
    uncorrelated): $V \to V/k$. The more independent the errors, the
    larger the gain. This is why we want diverse models (different samples,
    features, hyperparameters).
  \end{alertblock}
\end{frame}

%==================================================================
\section{Dropout}
%==================================================================

\begin{frame}[shrink=10]{Dropout: Motivation}
  \begin{block}{The Cost Problem with Ensembles}
    Model averaging (a bagging ensemble) almost always helps, but training
    several large neural networks for an ensemble is prohibitively
    expensive --- whether we train networks with different architectures
    (option 1) or multiple instances of the same network on different
    samples (option 2). Combining several models at \emph{test} time is
    also infeasible for real-time applications.
  \end{block}

  \begin{block}{What Dropout Does}
    Dropout addresses both issues. It effectively allows training of
    exponentially many neural networks without significant computational
    overhead, and gives an efficient approximate way of \emph{combining}
    them.
  \end{block}

  \begin{exampleblock}{Mechanism}
    Dropout means \textbf{dropping out units}: temporarily remove a node
    and its incoming/outgoing connections, producing a \emph{thinned}
    network. Each node is retained with a fixed probability $p$ (typically
    $p = 0.5$ for hidden nodes and $p = 0.8$ for visible/input nodes).
  \end{exampleblock}
\end{frame}

%------------------------------------------------------------------
\begin{frame}[shrink=10]{Dropout: Exponentially Many Thinned Networks}
  \begin{block}{Counting the Networks}
    With $n$ nodes, each can be retained or dropped, giving
    \[
      \boxed{\,2^{n}\,}
    \]
    possible thinned networks. This is prohibitively large --- we cannot
    train them all separately.
  \end{block}

  \begin{block}{The Two Tricks}
    \begin{enumerate}\setlength\itemsep{3pt}
      \item \textbf{Share the weights} across all thinned networks (one
            master weight matrix $W$).
      \item \textbf{Sample a different thinned network for each training
            instance} (or mini-batch).
    \end{enumerate}
  \end{block}

  \begin{exampleblock}{Training Loop}
    Initialise weights. For training instance / mini-batch $1$, apply
    dropout to get a thinned network, compute the loss, back-propagate ---
    updating \emph{only the active weights}. For instance $2$, apply dropout
    again to get a \emph{different} thinned network, and repeat. Because of
    weight sharing, the same master weights get updated through many
    different thinned networks over time.
  \end{exampleblock}
\end{frame}

%------------------------------------------------------------------
\begin{frame}[shrink=10]{Dropout: Parameter Sharing Across Networks}
  \begin{block}{Updates per Weight}
    \begin{itemize}\setlength\itemsep{3pt}
      \item If a weight was active in \emph{both} instances, by now it has
            received \emph{two} updates.
      \item If active in only \emph{one}, it has received \emph{one}
            update.
    \end{itemize}
    Each thinned network is trained rarely (or even never) on its own, but
    parameter sharing ensures \textbf{no weight is left untrained}: the
    shared $W$ is updated many times across the various thinned networks
    that contain it.
  \end{block}

  \begin{block}{What Happens at Test Time?}
    We have $2^{n}$ thinned networks; at test time each would have been used
    with probability $p$. It is impossible to aggregate the outputs of
    $2^{n}$ networks. Instead, we use the \textbf{full network} and
    \emph{scale} the output of each node by the fraction of time it was on
    during training (i.e.\ by $p$). A node always present at test time but
    on only a $p$-fraction of training time has its outgoing activation
    multiplied by $p$.
  \end{block}
\end{frame}

%------------------------------------------------------------------
\begin{frame}[shrink=10]{Dropout: Preventing Co-adaptation}
  \begin{block}{Masking Noise on Hidden Units}
    Dropout essentially applies a \textbf{masking noise} to the hidden
    units. It prevents hidden units from \emph{co-adapting}: a hidden unit
    cannot rely too much on other specific units, since any of them may be
    dropped at any time.
  \end{block}

  \begin{block}{Robustness Through Redundancy}
    Suppose hidden unit $h_i$ learns to detect a face by relying on another
    unit that detects the nose. If that nose-detector is dropped, $h_i$
    fails. Dropping a unit corresponds to \textbf{erasing the information}
    it encodes. To cope, the model must learn another unit that
    \emph{redundantly} encodes the nose, or learn to detect the face using
    \emph{other} features.
  \end{block}

  \begin{alertblock}{Intuition}
    Like discouraging lazy team members who depend on others: dropout
    forces every hidden unit to be useful on its own, so each learns to be
    more robust. ``$Koru\ ekule\ enkoru\ leru$'' --- no free-riding; each
    unit must pull its weight.
  \end{alertblock}
\end{frame}

%==================================================================
\section{Recap: Training Deep Neural Networks}
%==================================================================

\begin{frame}[shrink=10]{Gradient Descent and the Chain Rule}
  \begin{block}{Single Neuron}
    For a one-neuron network $\hat y = \sigma(wx)$ with loss $\mathcal L$,
    \[
      w = w - \eta\,\nabla w,
      \qquad
      \nabla w = \frac{\partial \mathcal L(w)}{\partial w}
        = \bigl(f(x)-y\bigr)\,f(x)\bigl(1-f(x)\bigr)\,x .
    \]
  \end{block}

  \begin{block}{Wider Network}
    With several inputs $x_1,x_2,x_3$ and weights $w_1,w_2,w_3$,
    \[
      w_j = w_j - \eta\,\frac{\partial \mathcal L}{\partial w_j},
      \qquad
      \nabla w_j = \bigl(f(x)-y\bigr)\,f(x)\bigl(1-f(x)\bigr)\,x_j .
    \]
  \end{block}

  \begin{block}{Deeper Network (Back-Propagation)}
    With $h_2 = \sigma(a_2)$, $a_2 = w_2 h_1$, $h_1 = \sigma(a_1)$,
    $a_1 = w_1 x = w_1 h_0$, the chain rule gives
    \[
      \nabla w_1 = \frac{\partial \mathcal L(w)}{\partial w_1}
        = \frac{\partial \mathcal L(w)}{\partial \hat y}\,
          \frac{\partial \hat y}{\partial a_2}\,
          \frac{\partial a_2}{\partial h_1}\,
          \frac{\partial h_1}{\partial a_1}\,
          \frac{\partial a_1}{\partial w_1}
        = (\cdots)\,h_0 .
    \]
  \end{block}
\end{frame}

%------------------------------------------------------------------
\begin{frame}[shrink=10]{The Two Key Observations}
  \begin{block}{Gradient $\propto$ Input}
    In general $\nabla w_i = \dfrac{\partial \mathcal L(w)}{\partial w_i}
    = (\cdots)\,h_{i-1}$: the gradient w.r.t.\ a parameter is
    \textbf{proportional to the input feeding into that parameter}
    ($h_{i-1}$ or, for the first layer, $x_i$). It directly depends on the
    activation of the preceding layer.
  \end{block}

  \begin{block}{Deep and Wide Networks}
    For a deep, wide network, $\nabla w_2$ is given by the chain rule
    applied across \textbf{multiple paths} --- this is exactly
    back-propagation. The gradients tell us the \emph{responsibility} of
    each parameter towards the loss, and they always have the same shape as
    the variable being differentiated.
  \end{block}

  \begin{exampleblock}{Historical Note}
    Back-propagation was popularised by Rumelhart \emph{et al.}\ in $1986$.
    It did not catch on immediately due to limited \textbf{computational
    power} and \textbf{memory}, and was not very successful for really deep
    networks. In fact, until $2006$ it was very hard to train very deep
    networks --- even after many epochs the training often did not converge.
  \end{exampleblock}
\end{frame}

%==================================================================
\section{DCT and JPEG Image Compression}
%==================================================================

\begin{frame}[shrink=10]{The Discrete Cosine Transform (DCT)}
  \begin{block}{Objective}
    Interpret the 2D DCT: represent an image as a \textbf{sum of basis
    images}. Compression then amounts to setting some DCT coefficients to
    zero --- this is the heart of \textbf{JPEG}.
  \end{block}

  \begin{block}{Forward 2D DCT}
    \[
      F(u,v)
        = \frac{2}{N}\,C(u)C(v)
          \sum_{x=0}^{N-1}\sum_{y=0}^{N-1}
          f(x,y)\,
          \cos\!\Bigl(\frac{(2x+1)u\pi}{2N}\Bigr)
          \cos\!\Bigl(\frac{(2y+1)v\pi}{2N}\Bigr)
    \]
    with $C(k) = \tfrac{1}{\sqrt2}$ if $k=0$ and $C(k)=1$ otherwise.
    Here $f(x,y)$ is the pixel value, $F(u,v)$ a DCT coefficient,
    $N$ the image size, $(x,y)$ spatial coordinates, $(u,v)$ frequency
    coordinates.
  \end{block}

  \begin{block}{Image as a Weighted Sum of Basis Images}
    \[
      X(m,n) = \sum_{k=0}^{N-1}\sum_{l=0}^{N-1}
        b_{k,l}\,C_{k,l}[m,n],
    \]
    where the weights $b_{k,l}$ are the discrete cosine coefficients and
    each basis image $C_{k,l}[m,n]$ is the product of a vertical frequency
    $\tfrac{k}{2N}$ cosine and a horizontal frequency $\tfrac{l}{2N}$
    cosine.
  \end{block}
\end{frame}

%------------------------------------------------------------------
\begin{frame}[shrink=10]{Worked DCT Example ($2\times2$ Image)}
  \begin{block}{Setup}
    Let $f(x,y) = \begin{bmatrix} 1 & 2\\ 3 & 4 \end{bmatrix}$, $N=2$,
    with $\cos(0)=1$, $\cos(\pi/4)=0.707$, $\cos(3\pi/4)=-0.707$.
  \end{block}

  \begin{block}{Compute $F(0,0)$ ($u=0,v=0$)}
    \[
      F(0,0) = \frac{2}{2}\,C(0)C(0)\sum_{x=0}^{1}\sum_{y=0}^{1} f(x,y)\cdot 1
        = \frac{1}{\sqrt2}\cdot\frac{1}{\sqrt2}\,(1+2+3+4)
        = \frac{1}{2}\cdot 10 = 5 .
    \]
    The DC coefficient $F(0,0)$ is proportional to the average pixel value.
  \end{block}

  \begin{block}{Other Coefficients}
    Carrying out the analogous sums gives, for this image,
    $F(0,1) \approx -1$, $F(1,0) \approx -2$, $F(1,1) = 0$, so the DCT
    matrix is
    \[
      F = \begin{bmatrix} 5 & -1\\ -2 & 0 \end{bmatrix}.
    \]
    A horizontal edge in the image becomes energy concentrated in the first
    column of the DCT matrix, because the image varies only in the vertical
    direction and has no horizontal frequency content.
  \end{block}
\end{frame}

%------------------------------------------------------------------
\begin{frame}[shrink=10]{Compression and the JPEG Pipeline}
  \begin{block}{Compression by Dropping Coefficients}
    Storing all coefficients reproduces the image exactly (no loss). Keeping
    only the few \emph{significant} coefficients (e.g.\ $4$ out of many)
    compresses with little visible loss, because most of the energy lives in
    a handful of low-frequency coefficients.
  \end{block}

  \begin{block}{JPEG Encoding Steps}
    \begin{enumerate}\setlength\itemsep{2pt}
      \item Colour-space conversion: $\text{RGB} \to Y C_b C_r$
            (luminance $Y$, blue chrominance $C_b$, red chrominance $C_r$).
            Reversible, no data compression.
      \item Chrominance down-sampling ($2\times2$ averaging of $C_b, C_r$):
            removes a considerable amount of data the eye barely notices.
      \item Break the $M\times N$ image into $8\times8$ blocks; take the 2D
            DCT of each block.
      \item \textbf{Quantisation}: divide each coefficient by a quantisation
            table and round to the nearest integer (this is where data is
            actually thrown away).
      \item Run-length encoding (RLE) and Huffman encoding of the
            quantised coefficients.
    \end{enumerate}
  \end{block}

  \begin{exampleblock}{JPEG Decoding}
    Read the bits for each $8\times8$ block, convert bits to DCT
    coefficients, take the inverse DCT to recover the $8\times8$ image, and
    put the blocks back together.
  \end{exampleblock}
\end{frame}

%------------------------------------------------------------------
\begin{frame}[shrink=10]{Why JPEG Works: Exploiting the Human Eye}
  \begin{block}{Eyes Are Not Perfect}
    The human eye has $\sim\!100$ million rod cells (brightness, low-light
    sensitive) but only $\sim\!6$ million cone cells (colour sensitive).
    So we perceive \textbf{brightness (luminance)} far better than
    \textbf{colour (chrominance)}. JPEG removes information we do not
    perceive well.
  \end{block}

  \begin{block}{DCT + Quantisation Together}
    The DCT alone does not compress; quantisation does. The
    \textbf{quantisation table} uses smaller divisors for low-frequency
    (top-left) coefficients --- which the eye notices --- and larger
    divisors for high-frequency (bottom-right) details the eye cannot
    perceive easily. After rounding, the high-frequency details are thrown
    away. We literally cannot tell the difference.
  \end{block}

  \begin{exampleblock}{High-Frequency Blindness}
    Eyes are good at seeing edges (e.g.\ the outline of a rock) but poor at
    high-frequency colour detail (single blades of grass, individual
    leaves in a cluster, shadow variation). The chrominance quantisation
    table uses higher numbers (coarser) than the luminance table. Then
    \textbf{run-length} encoding compactly stores the many resulting zeros
    (e.g.\ ``$0\times38$''), and \textbf{Huffman} coding compresses the
    rest. H.264 (AVC, Advanced Video Coding) extends these ideas to video.
  \end{exampleblock}
\end{frame}

%==================================================================
\section{Summary}
%==================================================================

\begin{frame}[shrink=12]{Summary: The Big Picture}
  \begin{columns}[T]
    \column{0.49\textwidth}
    \begin{block}{Autoencoders}
      Encode $\mathbf{h}=g(W\mathbf{x}+\mathbf{b})$, decode
      $\hat{\mathbf{x}}=f(W^{*}\mathbf{h}+\mathbf{c})$.\\
      Undercomplete: compresses (like PCA).\\
      Overcomplete: can learn identity $\Rightarrow$ needs regularisation.\\
      Binary input: sigmoid $+$ cross-entropy.\\
      Real input: linear $+$ squared error.
    \end{block}

    \begin{block}{PCA $\equiv$ Linear AE}
      With linear encoder/decoder, squared loss, and
      $\frac{1}{\sqrt m}$ centring: $H = X V_{\cdot,\le k}$, and the
      encoder $W = V_{\cdot,\le k}$ equals the PCA projection $P$
      (eigenvectors of $X^{\top}X$).
    \end{block}

    \column{0.49\textwidth}
    \begin{block}{Generalisation Theory}
      Expected error $=$ Bias$^2$ $+$ Variance $+$ noise.\\
      true error $=$ train error $+$ const $+\ \Omega(\text{complexity})$,\\
      $\Omega \approx \frac{\sigma^2}{n}\sum_i \partial\hat f(x_i)/\partial y_i$
      (Stein's lemma). Minimise
      $\mathcal L_{\text{train}}(\theta) + \Omega(\theta)$.
    \end{block}

    \begin{alertblock}{Regularisers}
      $L_2$ scales weights by $\frac{d_i}{d_i+\alpha}$; denoising; sparse
      (KL); contractive (Jacobian); weight tying; augmentation; input noise
      $\equiv L_2$; label smoothing; early stopping $\equiv L_2$; bagging
      ($V\!\to\!V/k$); dropout ($2^n$ shared networks).
    \end{alertblock}
  \end{columns}

  \medskip
  \begin{center}
    \large\textit{%
      ``An autoencoder learns a representation; regularisation stops it
        from cheating. Every regulariser is really one principle ---
        do not let the training error collapse to zero.''}
  \end{center}
\end{frame}

%------------------------------------------------------------------
\begin{frame}
  \Huge{\centerline{Thank You}}
  \vspace{0.6cm}
  \large{\centerline{Questions?}}
\end{frame}

\end{document}